% ============================================================
% f2014_sichuan_arts_q08_river_width.tex
% 2014 四川卷（文科）第 8 题：气球上测河流宽度
% 来源 PDF：普通高考/2014/2014四川文.pdf
% 原页：第 1 页（A3 横排）中栏，印面 921，第 8 题题干下方
% 说明：A 为气球位置，其正下方为河岸上的垂足；B、C 为对岸与
%       本岸的两个观测点，A 到 B、C 的俯角分别为 75°、30°；
%       气球高 60 m，图中给出水平虚线、铅垂高度线与直角标记。
% ============================================================
\begin{tikzpicture}[>=Stealth, line width=0.7pt, font=\small,
  x=1cm, y=1cm, line cap=round, line join=round]

  \coordinate (A) at (0,4);
  \coordinate (D) at (0,0);
  \coordinate (B) at (1.10,0);
  \coordinate (C) at (6.90,0);

  % 地面与过 A 的水平线
  \draw[dashed] (-0.80,0) -- (7.60,0);
  \draw[dashed] (-0.20,4) -- (2.20,4);

  % 视线
  \draw[dashed] (A) -- (B);
  \draw[dashed] (A) -- (C);

  % 气球
  \draw (0,4.44) circle (0.20);
  \draw (-0.08,4.26) -- (0.08,4.26) -- (0.10,4.00) -- (-0.10,4.00) -- cycle;
  \node at (0,4.45) {$A$};

  % 铅垂高度与标注
  \draw (-0.85,4) -- (-0.25,4);
  \draw (-0.85,0) -- (-0.25,0);
  \draw[<->] (-0.55,4) -- (-0.55,0);
  \node[left] at (-0.70,2) {$60\mathrm{m}$};

  % 垂足处的直角标记
  \draw (0,0.16) -- (0.16,0.16) -- (0.16,0);

  % 俯角
  \draw ($(A)+(0.60,0)$) arc (0:-30:0.60);
  \node at (0.98,3.80) {$30^\circ$};
  \draw ($(A)+(1.20,0)$) arc (0:-75:1.20);
  \node at (1.78,3.62) {$75^\circ$};

  % 点的标注
  \node[below] at (B) {$B$};
  \node[below] at (C) {$C$};
\end{tikzpicture}
